\section{成员选择与集成简化}
\label{sec:ensemble-selection-compression}

Bagging、Boosting和Stacking先回答怎样形成组合。训练完成后，成员数量仍可能超过存储或延迟预算，也可能包含只在少数输入区域有用的模型。简化集成有三种不同目标：为所有输入固定一个较小成员子集，针对当前输入临时选择成员，或训练一个新的紧凑模型替代原集成。前两种方法保留部分原成员，第三种方法改变了部署模型。

图\ref{fig:ensemble-selection-compression}沿部署对象区分这三种目标。静态选择把成员池缩成一个固定子集，动态选择保留成员池并让路由依赖当前输入，紧凑化则把完整集成的行为转移到新模型。后续评价资源收益时，需要测量图中实际部署的整条路径。

\begin{figure*}[htbp]
  \centering
  % Static selection, dynamic selection, and compression change deployment differently.
% Schematic only; no empirical quantities are encoded.
% Canvas: 164 mm x 64 mm. Dependencies: project palette and TikZ.
\begin{tikzpicture}[
  x=1mm,y=1mm,
  font=\footnotesize,
  text=BookInk,
  pool/.style={
    rectangle,
    draw=FigureModel,
    fill=FigureModel!8,
    line width=0.7pt,
    minimum width=42mm,
    minimum height=10mm,
    align=center,
    inner sep=1mm
  },
  process/.style={
    rectangle,
    draw=FigureLoss,
    fill=FigureLoss!9,
    line width=0.8pt,
    minimum width=42mm,
    minimum height=11mm,
    text width=38mm,
    align=center,
    inner sep=1mm
  },
  deployed/.style={
    rectangle,
    draw=FigureOutput,
    fill=FigureOutput!6,
    line width=0.8pt,
    minimum width=42mm,
    minimum height=12mm,
    text width=38mm,
    align=center,
    inner sep=1mm
  },
  flow/.style={
    -{Latex[length=2.2mm,width=1.5mm]},
    draw=BookMuted,
    line width=0.8pt
  }
]
  \path[use as bounding box] (0,0) rectangle (164,64);

  \node[font=\heiti\small] at (27,60) {(a) 静态成员选择};
  \node[pool] (static-pool) at (27,47)
    {已训练成员池\\$h_1,\ldots,h_M$};
  \node[process] (static-select) at (27,29)
    {依据统一验证目标\\固定子集$S$};
  \node[deployed] (static-deploy) at (27,10)
    {所有输入运行\\同一个成员子集};
  \draw[flow] (static-pool)--(static-select);
  \draw[flow] (static-select)--(static-deploy);

  \node[font=\heiti\small] at (82,60) {(b) 动态成员选择};
  \node[pool] (dynamic-pool) at (82,47)
    {已训练成员池\\$h_1,\ldots,h_M$};
  \node[process] (dynamic-select) at (82,29)
    {根据查询输入$\bm{x}$\\估计局部能力};
  \node[deployed] (dynamic-deploy) at (82,10)
    {不同输入可运行\\不同成员或权重};
  \draw[flow] (dynamic-pool)--(dynamic-select);
  \draw[flow] (dynamic-select)--(dynamic-deploy);

  \node[font=\heiti\small] at (137,60) {(c) 紧凑化};
  \node[pool] (teacher) at (137,47)
    {完整集成教师\\$H$};
  \node[process] (distill) at (137,29)
    {用标签与教师输出\\训练学生};
  \node[deployed] (student) at (137,10)
    {部署新的紧凑模型\\$s_{\bm{\theta}}$};
  \draw[flow] (teacher)--(distill);
  \draw[flow] (distill)--(student);

  \draw[BookRule,line width=0.6pt] (54.5,2)--(54.5,61);
  \draw[BookRule,line width=0.6pt] (109.5,2)--(109.5,61);
\end{tikzpicture}

  \caption{静态成员选择、动态成员选择与紧凑化改变的部署对象。静态选择为所有输入保留同一成员子集；动态选择根据输入改变运行成员或组合权重；紧凑化训练并部署新的学生模型。三种路径可能减少不同类型的成本，不能只用保留成员数比较。}
  \label{fig:ensemble-selection-compression}
\end{figure*}

\subsection{静态成员选择}
\label{sec:ensemble-static-selection}

\term{静态成员选择}（static ensemble selection）给定已经训练好的候选池$\{h_1,\ldots,h_M\}$，选择一个对所有未来输入都相同的子集$S\subseteq\{1,\ldots,M\}$。若$\mathcal{L}_{\mathrm{val}}(S)$表示该子集按指定规则组合后的验证损失，可以写成
\begin{equation}
  \widehat{S}
  \in
  \argmin_{S\subseteq\{1,\ldots,M\}}
  \mathcal{L}_{\mathrm{val}}(S)
  +\lambda_1|S|
  +\lambda_2 C(S),
  \label{eq:ensemble-static-selection}
\end{equation}
其中$\lambda_1,\lambda_2\geq0$，$|S|$控制成员数，$C(S)$可以表示存储、单样本延迟或其他部署成本。这个组合优化通常随成员池大小指数增长，需要用排序、前向加入、后向删除或其他近似搜索。

只按单成员验证损失排序，容易保留许多行为相近的模型；按当前组合的边际改进逐个加入成员，可以同时感知成员质量与已有组合。允许重复选择某个成员时，重复次数还形成离散权重。基于模型库的贪心集成选择就是这种做法的代表：从一个初始组合出发，每轮加入使验证指标改善最大的候选，必要时使用多起点或袋化减少搜索的不稳定性\scite{caruana2004ensemble}。

成员选择与一般模型选择的区别在于，式\eqref{eq:ensemble-static-selection}固定了候选池及其预测，只搜索哪些已训练成员共同部署。若同时改变基学习算法、特征处理和超参数，搜索范围已经扩展为完整模型开发流程。候选池越大，验证集上偶然出现有利组合的机会也越多；选择所用数据不能同时充当最终测试。

分歧率、误差协方差等多样性指标可以缩小候选范围，却不能替代任务损失。一个与其他成员高度不同但经常出错的模型，可能增加分歧而降低组合质量。静态选择应直接评价选定组合的输出，并在相同成员数或资源预算下与完整集成及单模型基线比较。

\subsection{动态成员选择}
\label{sec:ensemble-dynamic-selection}

静态选择假设同一成员子集适合所有输入。若成员在输入空间的不同区域各有所长，可以先估计当前输入附近的成员能力，再决定由谁预测。\term{动态分类器选择}（dynamic classifier selection, DCS）为每个输入选择一个成员；\term{动态集成选择}（dynamic ensemble selection, DES）选择一个成员子集或给出输入相关权重。

一种基本构造是准备独立的能力估计数据。对查询输入$\bm{x}$，在该数据中寻找非空邻域$\mathcal{N}(\bm{x})$，再以成员$m$在邻域内的正确率估计
\begin{equation}
  \widehat{c}_m(\bm{x})
  =
  \frac{1}{|\mathcal{N}(\bm{x})|}
  \sum_{i\in\mathcal{N}(\bm{x})}
  \mathbb{I}\!\left\{h_m(\bm{x}_i)=y_i\right\}.
  \label{eq:ensemble-local-competence}
\end{equation}
DCS可以选择$\widehat{c}_m(\bm{x})$最大的成员，DES则保留超过阈值的成员或按能力加权。局部正确率只是能力估计的一种方式，还可以使用输出轮廓、分类间隔或由元模型学习的能力特征\scite{woods1997localaccuracy}。

这个流程新增了两个误差来源。邻域依赖输入表示与距离，高维或分布稀疏时，附近样本未必提供稳定的局部证据；能力估计器也可能在有限数据上过拟合。成员训练数据、能力估计数据与最终测试数据需要分开，所有标准化和邻域参数都在相应开发流程内确定。

动态选择不一定降低端到端延迟。若系统必须先运行全部成员才能构造输出轮廓，再选择其中几个聚合，成员计算已经发生；基于输入邻域的方法还增加了检索成本。只有能够在运行昂贵成员之前完成路由，或动态组合确实改善资源分配时，成员数减少才可能转化为实际延迟收益。

\subsection{从集成模型到紧凑模型}
\label{sec:ensemble-compression}

静态与动态选择仍依赖原成员池。另一条路线是训练一个新的\term{紧凑模型}（compact model），使其逼近完整集成的输入输出映射。早期模型压缩工作已经用大型集成的输出训练较小模型；知识蒸馏随后把这种教师到学生的传递扩展为常用训练形式\scite{bucilua2006compression,hinton2015distilling}。

设完整集成为教师$H$，紧凑模型为学生$s_{\bm{\theta}}$。对回归，可以组合真实响应损失与教师拟合损失：
\begin{equation}
  \min_{\bm{\theta}}
  \sum_{i=1}^{n}
  \left[
    \alpha\ell\!\left(y_i,s_{\bm{\theta}}(\bm{x}_i)\right)
    +(1-\alpha)
    \ell\!\left(H(\bm{x}_i),s_{\bm{\theta}}(\bm{x}_i)\right)
  \right],
  \label{eq:ensemble-compression-objective}
\end{equation}
其中$\alpha\in[0,1]$控制真实标签与教师输出的作用。分类时可以拟合教师的完整类别分布；温度缩放会使非最大类别之间的相对关系更容易进入训练信号，但温度、损失权重与部署概率仍须分别选择和校准。

学生训练输入可以来自原训练集、额外未标注样本或经过说明的合成样本。教师只在这些输入上提供目标，输入覆盖不足时，学生即使在压缩数据上高度一致，也可能在其他区域偏离教师。学生容量过小会产生逼近误差，容量很大则可能无法达到预期资源节省。

紧凑化至少需要四类证据。预测一致性衡量学生是否复现教师输出；任务性能检查学生面对真实标签的风险；概率评价检查蒸馏后的得分是否仍可解释；资源测量确认模型大小、内存和延迟是否实际下降。教师与学生高度一致不自动保证二者具有相同新样本风险，学生也可能因真实标签项、正则化或额外数据而在部分样本上优于教师。原集成的成员级解释、袋外诊断和理论条件不能直接转移到新模型。
